\chapter{Metodologi}\label{ch:method}
% =============================================================================
% KERANGKA BAB 3 — isi tahapan dari PDF proposal revisi (Bab III, 13 tahap);
% susunan subbab mengikuti TA rujukan: 3.1 Data, 3.2 Variabel Penelitian,
% 3.3 Tahapan Penelitian (subsubbab per tahap), 3.4 Diagram Alir.
% Penyesuaian terhadap proposal yang wajib dicantumkan:
%  - D1 Opsi A: irisan 2.028 hari tanpa interpolasi; spline hanya sensitivitas (audit S2).
%  - Laju mlipir fase kedua dicantumkan di Tabel 3.2 (audit T5).
%  - Tahap verifikasi dan validasi mencakup batas bank latih dan batas dari bank
%    validasi (Tahap 21, D10 addendum).
% =============================================================================

    \section{Data}\label{sec:data source}
    % Data sekunder DJPPR Kementerian Keuangan: yield harian tujuh instrumen, kurs
    % USD/EUR dan USD/JPY, komposisi eksisting 2024, batas proporsi; 3 Juli 2017–
    % 24 November 2025 (2.028 hari perdagangan setelah irisan tanggal).
    % SUMBER: results/tables/audit_data.md

    \section{Variabel Penelitian}\label{sec:variables}

        \begin{table}[H]
            \centering
            \caption{Variabel penelitian}\label{tab:variables}
            % Isi: w_i (bobot), y_{i,t} (yield), X_{c,t} (kurs), C_{i,s} dan C^I_{i,s}
            % (biaya), eta dan tau (batas risiko) beserta keterangan dan satuan.
            \textit{[Tabel belum diisi]}
        \end{table}

    \section{Tahapan Penelitian}\label{sec:stages}

        \subsection{Studi Literatur}\label{subsec:literature study}
        \subsection{Pengumpulan Data}\label{subsec:data collection}
        \subsection{Pra-pemrosesan Data}\label{subsec:preprocessing}
        \subsection{Pengolahan Data}\label{subsec:processing}
        \subsection{Pembangkitan Skenario GBM}\label{subsec:scenario generation}
        \subsection{Formulasi Model MCCaR}\label{subsec:model formulation}
        \subsection{Skenario Batas Proporsi Portofolio Utang}\label{subsec:boundary scenarios}
        \subsection{Penentuan Batas Risiko}\label{subsec:risk limit determination}
        \subsection{Penyelesaian Optimasi}\label{subsec:solving}

            \begin{table}[H]
                \centering
                \caption{Parameter utama eksperimen}\label{tab:main parameters}
                % SUMBER: configs/params.yaml (blok risk, optimization, experiments, algorithms)
                % SUMBER: results/tables/tab_experiment_params.tex
                \textit{[Tabel belum diisi]}
            \end{table}

        \subsection{Hasil Optimasi}\label{subsec:optimization output}
        \subsection{Verifikasi dan Validasi}\label{subsec:verification}
        \subsection{Analisis Hasil dan Pembahasan}\label{subsec:analysis stage}
        \subsection{Penulisan Laporan Tugas Akhir}\label{subsec:writing}

    \section{Diagram Alir Penelitian}\label{sec:flowchart}

        \begin{figure}[H]
            \centering
            \textit{[Gambar belum disisipkan]}
            \caption{Diagram alir penelitian}\label{fig:flowchart}
            % Gambar 3.1 proposal; tambahkan cabang analisis tambahan (Tahap 16–21)
            % sebagai blok "Analisis Sensitivitas dan Validasi".
        \end{figure}
